\documentclass[11pt,letterpaper]{article}
\usepackage[margin=1in]{geometry}
\usepackage{amsmath}
\usepackage{amssymb}
\usepackage{parskip}

\title{The Geometric Distribution}
\author{}
\date{}

\begin{document}

\maketitle

\section*{What is the Geometric Distribution?}

The \textbf{geometric distribution} models the number of independent trials needed to achieve the first success in a sequence of Bernoulli trials (trials with only two outcomes: success with probability $p$ or failure with probability $1-p$).

\subsection*{Examples:}
\begin{itemize}
    \item Flipping a coin until you get heads
    \item Rolling a die until you get a 6
    \item Taking free throws until you make one
\end{itemize}

The probability mass function is:
\[
P(X = n) = (1-p)^{n-1} \cdot p
\]

where $n$ is the trial number of the first success.

\section*{Explaining the Last Step}

The simplification uses the \textbf{infinite geometric series formula}:

Starting with:
\[
\sum_{n=1}^{\infty} \binom{n-1}{0} 0.99^{n-1} \cdot 0.05 \cdot 0.95^{n-1}
\]

\textbf{Step 1:} Note that $\binom{n-1}{0} = 1$, so:
\[
= 0.05 \sum_{n=1}^{\infty} (0.99)^{n-1} \cdot (0.95)^{n-1}
\]

\textbf{Step 2:} Combine the terms:
\[
= 0.05 \sum_{n=1}^{\infty} (0.99 \cdot 0.95)^{n-1}
\]

\textbf{Step 3:} This is a geometric series $\sum_{k=0}^{\infty} r^k = \frac{1}{1-r}$ where $r = 0.99 \cdot 0.95$:
\[
= 0.05 \cdot \frac{1}{1 - 0.95 \cdot 0.99}
\]
\[
= \frac{0.05}{1 - 0.95 \cdot 0.99}
\]

The key insight is recognizing that when you have a sum of the form $\sum ar^n$, you can use the geometric series formula to get a closed-form expression.

\end{document}
